\documentclass[11pt]{article}
\usepackage[margin=1in]{geometry}
\usepackage{amsmath,amssymb,amsthm,hyperref}
\newtheorem{theorem}{Theorem}
\newcommand{\fl}[1]{\left\lfloor#1\right\rfloor}
\title{A squarefree floor sum and its class-number inversion}
\author{Principia Math In Progress}
\date{29 September 2026}
\begin{document}\maketitle
\begin{abstract}
We prove the odd-squarefree floor-sum identity of Chamberland and Dilcher
and its equivalent M\"obius-inversion formula. Expanding the number of
square roots modulo a squarefree integer into quadratic characters gives
an exact remainder sum; an elementary lattice count then recovers the
floor sum. The two equivalent source conjectures receive one campaign credit.
\end{abstract}

\section*{Task 1: the original conjecture}
For a positive integer $n$, define
\[
 f(n)=\sum_{j=1}^{\lfloor n/4\rfloor}\fl{\sqrt{jn}}-\frac{n^2-1}{12}.
\]
For odd squarefree $d$, let $\delta(d)=1$ when $d\equiv3\pmod4$ and
$\delta(d)=0$ otherwise. Define $H(d)=0$ in the latter case, and
\[
 H(3)=\frac13,\qquad
 H(d)=h\bigl(\mathbb Q(\sqrt{-d})\bigr)
       \quad(d>3,\ d\equiv3\pmod4).
\]
Thus $H(d)$ is the source's product $h^*(-d)\delta(d)$, with no need
to assign an unused normalization when $\delta(d)=0$.

\begin{theorem}\label{thm:main}
For every odd squarefree positive integer $n$,
\begin{align}
 f(n)&=\frac{1-n}{4}\delta(n)-\frac12\sum_{d\mid n}H(d),
 \label{eq:floor}\\
 H(n)&=\sum_{d\mid n}\mu(n/d)
       \left(\frac{1-d}{2}\delta(d)-2f(d)\right).
 \label{eq:inversion}
\end{align}
\end{theorem}
Equation~\eqref{eq:inversion} is Conjecture 7.4 of~\cite[p.~22]{CD},
MathDB 370788. Equation~\eqref{eq:floor} is Conjecture 7.3,
MathDB 370787. M\"obius inversion makes them equivalent; we count them
as one resolution. The case $n=1$ has only empty sums and is immediate.

\section*{Counting square roots and evaluating the remainder sum}
For odd squarefree $d>1$, write $\chi_d(t)=(t/d)$ for the Jacobi symbol;
put $\chi_1(t)=1$. This is a nonprincipal quadratic character modulo $d$
when $d>1$. Its character sum is zero. One elementary way to see this is
to choose, by the Chinese remainder theorem, a unit that is a quadratic
nonresidue at exactly one prime dividing $d$; multiplication by this
unit negates the sum.

The number $\rho_n(t)$ of solutions of $x^2\equiv t\pmod n$ is
\begin{equation}\label{eq:roots}
 \rho_n(t)=\prod_{p\mid n}\left(1+\left(\frac tp\right)\right)
         =\sum_{d\mid n}\chi_d(t).
\end{equation}
This includes nonunit $t$: modulo a prime there is one root at zero
and zero or two roots at a nonzero residue. The Chinese remainder
theorem then proves the product for every $t$.

Define
\[
 S_n=\frac1n\sum_{x=0}^{n-1}(x^2\bmod n).
\]
Interchanging finite sums in~\eqref{eq:roots} gives
\[
 S_n=\sum_{d\mid n}\frac1n\sum_{t=0}^{n-1}t\chi_d(t).
\]
The contribution of $d=1$ is $(n-1)/2$. For $d>1$, split $t$ into
complete periods modulo $d$. Since the unweighted character sum is zero,
\[
 \frac1n\sum_{t=0}^{n-1}t\chi_d(t)
 =\frac1d\sum_{r=0}^{d-1}r\chi_d(r).
\]
If $d\equiv1\pmod4$, the character is even, and pairing $r$ with
$d-r$ shows that this weighted sum is zero. If $d\equiv3\pmod4$,
quadratic reciprocity identifies $\chi_d(r)$ with the Kronecker
character $(-d/r)$. The discriminant $-d$ is fundamental, so
Dirichlet's finite class-number formula gives
\[
 \frac1d\sum_{r=0}^{d-1}r\chi_d(r)=-H(d).
\]
Here we use the classical formula stated in~\cite[Theorem 4.1]{CD};
its unit factor is six at $d=3$ and two otherwise, giving exactly
$H(3)=1/3$. Consequently
\begin{equation}\label{eq:remainder}
 S_n=\frac{n-1}{2}-\sum_{d\mid n}H(d).
\end{equation}

\section*{Recovering the floor sum}
We include the lattice-count step to fix every endpoint. Put
$\nu=\lfloor n/4\rfloor$ and
$F(n)=f(n)+(n^2-1)/12$. For odd $n\ge3$ one has
$\lfloor\sqrt{\nu n}\rfloor=2\nu$, including $\nu=0$.
Counting the pairs $(j,k)$ with $1\le j\le\nu$ and $k^2\le jn$ gives
\[
 F(n)=\sum_{k=1}^{2\nu}
       \left(\nu+1-\left\lceil\frac{k^2}{n}\right\rceil\right)
     =2\nu^2-\sum_{k=1}^{2\nu}\fl{\frac{k^2}{n}}.
\]
The last equality uses squarefreeness: $n\nmid k^2$ for $0<k<n$.
Writing $R_n(k)=k^2\bmod n$ and summing the consecutive squares yields
\begin{equation}\label{eq:lattice}
 F(n)=2\nu^2-\frac{2\nu(2\nu+1)(4\nu+1)}{6n}
             +\frac1n\sum_{k=1}^{2\nu}R_n(k).
\end{equation}
Since $R_n(n-k)=R_n(k)$,
\[
 \frac1n\sum_{k=1}^{(n-1)/2}R_n(k)=\frac{S_n}{2}.
\]
If $n\equiv1\pmod4$, then $2\nu=(n-1)/2$. If $n\equiv3\pmod4$,
then $2\nu=(n-3)/2$, and the omitted last remainder is
\[
 R_n\left(\frac{n-1}{2}\right)=\frac{n+1}{4},
\]
as follows by subtracting $(n+1)/4$ from $(n-1)^2/4$.
Substituting these two cases into~\eqref{eq:lattice} and subtracting
$(n^2-1)/12$ gives respectively
\[
 f(n)=\frac{S_n}{2}-\frac{n-1}{4},
 \qquad
 f(n)=\frac{S_n}{2}-\frac{n-1}{2}.
\]
Equation~\eqref{eq:remainder} now proves~\eqref{eq:floor}.

Finally, multiply~\eqref{eq:floor} by $-2$ and rearrange:
\[
 \frac{1-n}{2}\delta(n)-2f(n)=\sum_{d\mid n}H(d).
\]
Every divisor of an odd squarefree integer is again odd and squarefree.
Ordinary M\"obius inversion on the divisor set therefore gives
Equation~\eqref{eq:inversion}, completing the proof.

\section*{Tasks 2--5: follow-up questions}
\begin{enumerate}\setcounter{enumi}{1}
\item Extend the root-count expansion to arbitrary odd nonsquarefree
integers and obtain explicit coefficients for all class-number terms.
\item Derive effective uniform bounds for the deviation of $f(n)$ from
its residue-class main term over odd squarefree integers.
\item Replace the linear remainder weight by higher powers and express
the resulting sums using generalized Bernoulli numbers or Dirichlet values.
\item Compare algorithms for computing the class numbers in
Equation~\eqref{eq:inversion}, including their exact arithmetic complexity,
against enumeration of reduced binary quadratic forms.
\end{enumerate}

\paragraph{Verification and scope.}
The verifier computes the original floor and remainder sums using integer
arithmetic and computes class numbers independently by enumerating reduced
binary quadratic forms. It checks both identities and the inversion.
Finite checks corroborate the all-squarefree proof. This is author-agent
review, not independent peer review. Source and search limitations are
recorded in the literature audit.

\begin{thebibliography}{9}
\bibitem{CD} M.~Chamberland and K.~Dilcher,
\emph{Sums of the floor function related to class numbers of imaginary
quadratic fields},
\href{https://arxiv.org/abs/2510.04387}{arXiv:2510.04387v1} (2025),
Conjectures 7.3--7.4, pp.~21--22.
All 23 pages of the current manuscript were read on 29 September 2026.
\bibitem{MathDB} MathDB,
\href{https://mathdb.com/p/370788}{problem 370788} and the equivalent
\href{https://mathdb.com/p/370787}{problem 370787},
live records checked 29 September 2026.
\end{thebibliography}
\end{document}
